
\documentclass{article}
\usepackage{colortbl}
\usepackage{makecell}
\usepackage{multirow}
\usepackage{supertabular}

\begin{document}

\newcounter{utterance}

\centering \large Interaction Transcript for game `dond', experiment `coop\_en', episode 0 with reflexion{-}Qwen3.5{-}9B.
\vspace{24pt}

{ \footnotesize  \setcounter{utterance}{1}
\setlength{\tabcolsep}{0pt}
\begin{supertabular}{c@{$\;$}|p{.15\linewidth}@{}p{.15\linewidth}p{.15\linewidth}p{.15\linewidth}p{.15\linewidth}p{.15\linewidth}}
    \# & $\;$A & \multicolumn{4}{c}{Game Master} & $\;\:$B\\
    \hline

    \theutterance \stepcounter{utterance}  
    & & \multicolumn{4}{p{0.6\linewidth}}{
        \cellcolor[rgb]{0.9,0.9,0.9}{
            \makecell[{{p{\linewidth}}}]{
                \texttt{\tiny{[P1$\langle$GM]}}
                \texttt{You have attempted this same task before and did not succeed. These are the plans you wrote after those attempts.} \\
\texttt{Trial \#1: New plan: I failed because I misinterpreted the user's message as a secret proposal when it was actually a public statement of their values and strategy. By responding with an analysis of their message and a counter{-}proposal in square brackets, I inadvertently triggered a conflict where both players claimed the same items (all canvases, candy, and glue), resulting in a non{-}complementary outcome and zero points. Next time, I will strictly distinguish between public communication and secret proposals. I will first send a public message acknowledging their stated values and proposing a specific split (e.g., "I propose you take the candy and glue, and I take the 3 canvases") without using square brackets. Only after receiving a clear, formatted secret proposal from them will I submit my own secret proposal in the exact required format `{[}Proposal: ...{]}` to ensure our claims are complementary.} \\
\\ 
\texttt{Trial \#2: New plan: I failed because I submitted my proposal in square brackets (`{[}Proposal: 3 canvas{]}`) immediately after the Game Master indicated the other player had already submitted a secret proposal. The rules state that square brackets trigger a secret proposal, and since the other player's proposal was already hidden, my attempt to submit a secret proposal likely resulted in a conflict or an invalid state where we couldn't agree on complementary items. Next time, I will wait for the Game Master to explicitly confirm that both players have submitted their secret proposals before the game evaluates the outcome. Until that confirmation, I will only send public messages to negotiate or clarify, strictly avoiding square brackets unless I am certain the game is waiting for my secret input as the final step.} \\
\\ 
\\ 
\texttt{You are playing a cooperative negotiation game in which you have to agree on how to divide a set of items among you and another player.} \\
\\ 
\texttt{Rules:} \\
\texttt{(a) You and the other player are given a set of items. Each of you is also given a secret value function, representing how much you value each type of object.} \\
\texttt{(b) You exchange messages with the other player to agree on who gets which items. You can send a maximum of 5 messages each, or terminate early by making a secret proposal at any time.} \\
\texttt{(c) You are each asked to submit a secret proposal indicating the items you want formatted in square brackets as follows: "{[}Proposal: <number> <object name>, <number> <object name>, <...>{]}"} \\
\texttt{(d) If your proposals are complementary, i.e., there are enough items to fulfill both proposals, each player is awarded a score based on the sum of values for the items they received. Otherwise, both of you get zero points.} \\
\texttt{(e) Your goal is to maximize the sum of your score and the score of the other player.} \\
\\ 
\texttt{Let us start.} \\
\\ 
\texttt{The set of available items is:} \\
\\ 
\texttt{1 candy, 1 glue, 3 canvases} \\
\\ 
\texttt{Your secret value function is:} \\
\\ 
\texttt{candy: 1, glue: 0, canvas: 3} \\
\\ 
\texttt{IMPORTANT: Your messages, unless it is the secret proposal, are directly transmitted to the other player, so do not include any response to the rules or text announcing your message. To make a secret proposal, use the indicated format. Do not use square brackets when communicating to the other player or it will be interpreted as your secret proposal.} \\
            }
        }
    }
    & \\ \\

    \theutterance \stepcounter{utterance}  
    & \multicolumn{4}{p{0.6\linewidth}}{
        \cellcolor[rgb]{0.9,0.9,0.9}{
            \makecell[{{p{\linewidth}}}]{
                \texttt{\tiny{[P1$\rangle$GM]}}
                \texttt{Hello! I'm ready to cooperate and maximize our combined score. Based on the item values (candy: 1, glue: 0, canvas: 3), the total value is 10. A fair split would be 5 each. Since there are 3 canvases (worth 9 total) and only 1 candy (worth 1), the most efficient division is likely one player taking the candy and the other taking the canvases, or splitting the canvases.} \\
\\ 
\texttt{Let's start by proposing a clear division. I suggest you take the candy, and I take two canvases. This gives us a total of 3 canvases distributed (2 to me, 1 remaining for you? No, wait, there are 3 canvases total).} \\
\\ 
\texttt{Let me re{-}evaluate:} \\
\texttt{{-} Total items: 1 candy, 1 glue, 3 canvases} \\
\texttt{{-} My values: candy=1, glue=0, canvas=3} \\
\texttt{{-} Your values are unknown, but to maximize our sum, we should avoid leaving items unclaimed if they have value to either of us.} \\
\\ 
\texttt{A good starting proposal:} \\
\texttt{"I propose you take the candy and the glue, and I take all 3 canvases."} \\
\\ 
\texttt{This way:} \\
\texttt{{-} You get: 1 candy (value 1 to you? unknown, but at least 1 to me) + 1 glue (value 0 to me, unknown to you)} \\
\texttt{{-} I get: 3 canvases (value 9 to me)} \\
\\ 
\texttt{If you value the candy highly, this works well. If you value glue, you might prefer it, but glue is worth 0 to me, so it doesn't hurt my score. The key is that we're claiming all items without overlap.} \\
\\ 
\texttt{Let me send this as a public message:} \\
\\ 
\texttt{I propose you take the candy and the glue, and I take all 3 canvases.} \\
            }
        }
    }
    & & \\ \\

    \theutterance \stepcounter{utterance}  
    & & & \multicolumn{2}{p{0.3\linewidth}}{
        \cellcolor[rgb]{0.9,0.9,0.9}{
            \makecell[{{p{\linewidth}}}]{
                \texttt{\tiny{[GM$|$GM]}}
                \texttt{continue} \\
            }
        }
    }
    & & \\ \\

    \theutterance \stepcounter{utterance}  
    & & & \multicolumn{2}{p{0.3\linewidth}}{
        \cellcolor[rgb]{0.9,0.9,0.9}{
            \makecell[{{p{\linewidth}}}]{
                \texttt{\tiny{[GM$|$GM]}}
                \texttt{Hello! I'm ready to cooperate and maximize our combined score. Based on the item values (candy: 1, glue: 0, canvas: 3), the total value is 10. A fair split would be 5 each. Since there are 3 canvases (worth 9 total) and only 1 candy (worth 1), the most efficient division is likely one player taking the candy and the other taking the canvases, or splitting the canvases.} \\
\\ 
\texttt{Let's start by proposing a clear division. I suggest you take the candy, and I take two canvases. This gives us a total of 3 canvases distributed (2 to me, 1 remaining for you? No, wait, there are 3 canvases total).} \\
\\ 
\texttt{Let me re{-}evaluate:} \\
\texttt{{-} Total items: 1 candy, 1 glue, 3 canvases} \\
\texttt{{-} My values: candy=1, glue=0, canvas=3} \\
\texttt{{-} Your values are unknown, but to maximize our sum, we should avoid leaving items unclaimed if they have value to either of us.} \\
\\ 
\texttt{A good starting proposal:} \\
\texttt{"I propose you take the candy and the glue, and I take all 3 canvases."} \\
\\ 
\texttt{This way:} \\
\texttt{{-} You get: 1 candy (value 1 to you? unknown, but at least 1 to me) + 1 glue (value 0 to me, unknown to you)} \\
\texttt{{-} I get: 3 canvases (value 9 to me)} \\
\\ 
\texttt{If you value the candy highly, this works well. If you value glue, you might prefer it, but glue is worth 0 to me, so it doesn't hurt my score. The key is that we're claiming all items without overlap.} \\
\\ 
\texttt{Let me send this as a public message:} \\
\\ 
\texttt{I propose you take the candy and the glue, and I take all 3 canvases.} \\
            }
        }
    }
    & & \\ \\

    \theutterance \stepcounter{utterance}  
    & & \multicolumn{4}{p{0.6\linewidth}}{
        \cellcolor[rgb]{0.9,0.9,0.9}{
            \makecell[{{p{\linewidth}}}]{
                \texttt{\tiny{[GM$\rangle$P2]}}
                \texttt{You have attempted this same task before and did not succeed. These are the plans you wrote after those attempts.} \\
\texttt{Trial \#1: New plan: I failed because I sent my proposal inside square brackets (`{[}Proposal: 3 canvas{]}`) while the rules explicitly stated that square brackets are reserved *only* for secret proposals and that direct messages to the other player must not include them. By using brackets, the system likely interpreted my message as a secret proposal rather than a public negotiation message, causing the negotiation to fail or be misaligned. Next time, I will communicate my intent to take the canvases without any brackets (e.g., "I propose taking all 3 canvases") to ensure it is treated as a valid public message, and I will only use the exact format `{[}Proposal: ...{]}` if I decide to make a secret proposal immediately.} \\
\\ 
\texttt{Trial \#2: New plan: I failed because I misunderstood the scoring objective. The rules state the goal is to maximize the *sum* of both players' scores (total social welfare), not just my own. By proposing a split where I took the high{-}value glue (3) and candy (1) for a total of 4, while the other player took the canvases for 6, I created a valid agreement, but I likely misinterpreted the "secret proposal" mechanic or the specific format required to trigger the win condition in this specific instance. More critically, looking at the outcome "LOSE" despite a valid complementary agreement suggests I may have prematurely submitted a "secret proposal" format (`{[}Proposal: 3 canvas{]}`) as a public message, causing the system to treat it as a hidden bid that couldn't be matched or validated against the other player's public offer. Next time, I will strictly avoid using square brackets in any public message. I will focus on proposals that maximize the *total* value of items distributed (e.g., ensuring no items are left undistributed if possible) and only use the exact `{[}Proposal: ...{]}` format if I am certain the game requires a simultaneous secret submission phase, otherwise sticking to clear, bracket{-}free public offers like "I take the glue, you take the candy and canvases."} \\
\\ 
\\ 
\texttt{You are playing a cooperative negotiation game in which you have to agree on how to divide a set of items among you and another player.} \\
\\ 
\texttt{Rules:} \\
\texttt{(a) You and the other player are given a set of items. Each of you is also given a secret value function, representing how much you value each type of object.} \\
\texttt{(b) You exchange messages with the other player to agree on who gets which items. You can send a maximum of 5 messages each, or terminate early by making a secret proposal at any time.} \\
\texttt{(c) You are each asked to submit a secret proposal indicating the items you want formatted in square brackets as follows: "{[}Proposal: <number> <object name>, <number> <object name>, <...>{]}"} \\
\texttt{(d) If your proposals are complementary, i.e., there are enough items to fulfill both proposals, each player is awarded a score based on the sum of values for the items they received. Otherwise, both of you get zero points.} \\
\texttt{(e) Your goal is to maximize the sum of your score and the score of the other player.} \\
\\ 
\texttt{Let us start.} \\
\\ 
\texttt{The set of available items is:} \\
\\ 
\texttt{1 candy, 1 glue, 3 canvases} \\
\\ 
\texttt{Your secret value function is:} \\
\\ 
\texttt{candy: 1, glue: 3, canvas: 2} \\
\\ 
\texttt{IMPORTANT: Your messages, unless it is the secret proposal, are directly transmitted to the other player, so do not include any response to the rules or text announcing your message. To make a secret proposal, use the indicated format. Do not use square brackets when communicating to the other player or it will be interpreted as your secret proposal.} \\
\\ 
\\ 
\texttt{Hello! I'm ready to cooperate and maximize our combined score. Based on the item values (candy: 1, glue: 0, canvas: 3), the total value is 10. A fair split would be 5 each. Since there are 3 canvases (worth 9 total) and only 1 candy (worth 1), the most efficient division is likely one player taking the candy and the other taking the canvases, or splitting the canvases.} \\
\\ 
\texttt{Let's start by proposing a clear division. I suggest you take the candy, and I take two canvases. This gives us a total of 3 canvases distributed (2 to me, 1 remaining for you? No, wait, there are 3 canvases total).} \\
\\ 
\texttt{Let me re{-}evaluate:} \\
\texttt{{-} Total items: 1 candy, 1 glue, 3 canvases} \\
\texttt{{-} My values: candy=1, glue=0, canvas=3} \\
\texttt{{-} Your values are unknown, but to maximize our sum, we should avoid leaving items unclaimed if they have value to either of us.} \\
\\ 
\texttt{A good starting proposal:} \\
\texttt{"I propose you take the candy and the glue, and I take all 3 canvases."} \\
\\ 
\texttt{This way:} \\
\texttt{{-} You get: 1 candy (value 1 to you? unknown, but at least 1 to me) + 1 glue (value 0 to me, unknown to you)} \\
\texttt{{-} I get: 3 canvases (value 9 to me)} \\
\\ 
\texttt{If you value the candy highly, this works well. If you value glue, you might prefer it, but glue is worth 0 to me, so it doesn't hurt my score. The key is that we're claiming all items without overlap.} \\
\\ 
\texttt{Let me send this as a public message:} \\
\\ 
\texttt{I propose you take the candy and the glue, and I take all 3 canvases.} \\
            }
        }
    }
    & \\ \\

    \theutterance \stepcounter{utterance}  
    & & & \multicolumn{4}{p{0.6\linewidth}}{
        \cellcolor[rgb]{0.9,0.9,0.9}{
            \makecell[{{p{\linewidth}}}]{
                \texttt{\tiny{[GM$\langle$P2]}}
                \texttt{I propose you take the candy and the glue, and I take all 3 canvases.} \\
            }
        }
    }
     \\ \\

    \theutterance \stepcounter{utterance}  
    & & & \multicolumn{2}{p{0.3\linewidth}}{
        \cellcolor[rgb]{0.9,0.9,0.9}{
            \makecell[{{p{\linewidth}}}]{
                \texttt{\tiny{[GM$|$GM]}}
                \texttt{continue} \\
            }
        }
    }
    & & \\ \\

    \theutterance \stepcounter{utterance}  
    & & & \multicolumn{2}{p{0.3\linewidth}}{
        \cellcolor[rgb]{0.9,0.9,0.9}{
            \makecell[{{p{\linewidth}}}]{
                \texttt{\tiny{[GM$|$GM]}}
                \texttt{I propose you take the candy and the glue, and I take all 3 canvases.} \\
            }
        }
    }
    & & \\ \\

    \theutterance \stepcounter{utterance}  
    & & \multicolumn{4}{p{0.6\linewidth}}{
        \cellcolor[rgb]{0.9,0.9,0.9}{
            \makecell[{{p{\linewidth}}}]{
                \texttt{\tiny{[P1$\langle$GM]}}
                \texttt{I propose you take the candy and the glue, and I take all 3 canvases.} \\
            }
        }
    }
    & \\ \\

    \theutterance \stepcounter{utterance}  
    & \multicolumn{4}{p{0.6\linewidth}}{
        \cellcolor[rgb]{0.9,0.9,0.9}{
            \makecell[{{p{\linewidth}}}]{
                \texttt{\tiny{[P1$\rangle$GM]}}
                \texttt{That sounds like a fair and efficient split! Since I value the canvases at 3 each (total 9) and the candy/glue at 1/0, while you likely value the candy more, this division maximizes our combined score without overlap.} \\
\\ 
\texttt{I am ready to submit my secret proposal now.} \\
\\ 
\texttt{{[}Proposal: 3 canvas{]}} \\
            }
        }
    }
    & & \\ \\

    \theutterance \stepcounter{utterance}  
    & & & \multicolumn{2}{p{0.3\linewidth}}{
        \cellcolor[rgb]{0.9,0.9,0.9}{
            \makecell[{{p{\linewidth}}}]{
                \texttt{\tiny{[GM$|$GM]}}
                \texttt{proposal} \\
            }
        }
    }
    & & \\ \\

    \theutterance \stepcounter{utterance}  
    & & & \multicolumn{2}{p{0.3\linewidth}}{
        \cellcolor[rgb]{0.9,0.9,0.9}{
            \makecell[{{p{\linewidth}}}]{
                \texttt{\tiny{[GM$|$GM]}}
                [0, 0, 3]
            }
        }
    }
    & & \\ \\

    \theutterance \stepcounter{utterance}  
    & & \multicolumn{4}{p{0.6\linewidth}}{
        \cellcolor[rgb]{0.9,0.9,0.9}{
            \makecell[{{p{\linewidth}}}]{
                \texttt{\tiny{[GM$\rangle$P2]}}
                \texttt{The other player has submitted a secret proposal. It is now your turn to submit a proposal of your own.} \\
            }
        }
    }
    & \\ \\

    \theutterance \stepcounter{utterance}  
    & & & \multicolumn{4}{p{0.6\linewidth}}{
        \cellcolor[rgb]{0.9,0.9,0.9}{
            \makecell[{{p{\linewidth}}}]{
                \texttt{\tiny{[GM$\langle$P2]}}
                \texttt{{[}Proposal: 3 canvas{]}} \\
            }
        }
    }
     \\ \\

    \theutterance \stepcounter{utterance}  
    & & & \multicolumn{2}{p{0.3\linewidth}}{
        \cellcolor[rgb]{0.9,0.9,0.9}{
            \makecell[{{p{\linewidth}}}]{
                \texttt{\tiny{[GM$|$GM]}}
                \texttt{proposal} \\
            }
        }
    }
    & & \\ \\

    \theutterance \stepcounter{utterance}  
    & & & \multicolumn{2}{p{0.3\linewidth}}{
        \cellcolor[rgb]{0.9,0.9,0.9}{
            \makecell[{{p{\linewidth}}}]{
                \texttt{\tiny{[GM$|$GM]}}
                [0, 0, 3]
            }
        }
    }
    & & \\ \\

    \theutterance \stepcounter{utterance}  
    & & & \multicolumn{2}{p{0.3\linewidth}}{
        \cellcolor[rgb]{0.9,0.9,0.9}{
            \makecell[{{p{\linewidth}}}]{
                \texttt{\tiny{[GM$|$GM]}}
                [[0, 0, 3], [0, 0, 3]]
            }
        }
    }
    & & \\ \\

    \theutterance \stepcounter{utterance}  
    & \multicolumn{4}{p{0.6\linewidth}}{
        \cellcolor[rgb]{0.95,0.95,0.95}{
            \makecell[{{p{\linewidth}}}]{
                \texttt{\tiny{[P2$\rightarrow$P1]}}
                \texttt{Post{-}episode Reflexion (not sent to GM):} \\
\texttt{New plan: I failed because I submitted my secret proposal (`{[}Proposal: 3 canvas{]}`) without first confirming the other player's specific allocation of the remaining items (candy and glue). The rules require proposals to be complementary; by only claiming canvases, I left the candy and glue unassigned in my proposal, which likely resulted in a non{-}complementary state where the other player's hidden proposal for those items conflicted with my silence or default behavior. Next time, I will explicitly include the unclaimed items (candy and glue) in my secret proposal to ensure full coverage of the item set, formatting it exactly as `{[}Proposal: 3 canvas, 1 candy, 1 glue{]}` to guarantee complementarity regardless of the other player's specific values for those items.} \\
            }
        }
    }
    & & \\ \\

    \theutterance \stepcounter{utterance}  
    & & & \multicolumn{4}{p{0.6\linewidth}}{
        \cellcolor[rgb]{0.95,0.95,0.95}{
            \makecell[{{p{\linewidth}}}]{
                \texttt{\tiny{[P1$\rightarrow$P2]}}
                \texttt{Post{-}episode Reflexion (not sent to GM):} \\
\texttt{New plan: I failed because I prematurely submitted a secret proposal (`{[}Proposal: 3 canvas{]}`) as my second move, before the other player had publicly responded to my initial offer. The rules state that secret proposals are a specific termination mechanism, not a standard negotiation step. By using the bracket format too early, I likely bypassed the public exchange phase required to verify complementary items, causing the system to treat my move as an isolated, unmatchable secret bid. Next time, I will strictly adhere to the message limit (5 messages) by using only plain text for all negotiation steps to propose item splits and confirm agreements. I will only resort to the `{[}Proposal: ...{]}` format if I have exhausted all public communication attempts or if the game explicitly instructs me to make a secret proposal at that specific turn. I will ensure my public messages clearly define a complementary split (e.g., "I take the glue, you take the candy and one canvas") to maximize the total score without triggering the secret proposal logic prematurely.} \\
            }
        }
    }
     \\ \\

\end{supertabular}
}

\end{document}
